% 원제: 17. 결론
\chapter{결론}

PRISM-Hub는 고정비 대전력 경로와 부분전력 조정을 결합해, 800 V급 차량의 충전 호환과 다중 전압 보조전원을 하나의 직류 전력 구조로 다루는 아키텍처다. 현재까지의 검증으로 연구 질문의 답이 다음과 같이 정리되었다.

\begin{itemize}
\item \textbf{Q1(정격 분리): 긍정적.} 주행과 150 kW 충전 중 모두 링크전력이 8 kW 안이고, 충전전류는 조정단을 지나지 않는다. R2에서는 충전기가 Stage A 출력 노드에 직접 붙으므로, 충전 중 Stage A 정지를 충전기 차단과 협조시키는 V\textsubscript{A} 클램프와 정지 순서를 함께 정했다.
\item \textbf{Q2(강건 제어와 부하 관리): 긍정적.} 같은 제어기가 기존 구성과 창 조정 모듈 모두에서 공차·지연을 포함한 기준을 지켰다. 고장 조건에서는 측정 기반 상태기계가 Stage A를 멈추지 않고 Stage B를 우회하며, 실제 Stage A·조정형 Stage C를 넣은 통합 회로에서도 강등 중 48 V가 ±0.03 \% 안이었다. 회로 확인에서 드러난 현상(DAB 순환전류, 데드타임, 지령 전이, Stage C 영전압 전이, 역방향 경부하 이득, Stage A 경부하 영전류 스위칭)은 모두 검증 모델로 재현한 뒤 설계를 고쳤다. 영교차 게이트의 방향 반전 1건이 남았다.
\item \textbf{Q3(기준 구조 대비 이점): 조건부.} 전 범위 400 V 조정은 손실·정격 면에서 이점이 없어 H5가 기각되었다. 팩 등급 조건부 창 조정 모듈은 좁은 범위에서 성립 구조 가운데 손실이 가장 작다. 산업 방식과 비교하면 R2의 이점은 400 V급 보조부품을 재사용하는 조건에서만 성립한다. 이 조건에서 모터·인버터 승압 구조보다 보조 경로 손실이 26 \% 작고 충전 손실도 작지만, 800 V급 보조부품 설계보다는 보조 손실이 크다. 구성 원리는 선행이 있고, 기여는 그 조합과 성립 조건, 검증으로 얻은 설계 규칙이다.
\end{itemize}

이에 따라 아키텍처를 “Stage A 중심 다기능 허브 + 팩 등급 조건부 창 조정 모듈(R2)”로 정하고, 기본형(Stage A + 조정형 Stage C)까지 통합 회로에서 확인했다. 부정적 결과(H5 기각, 800 V급 부품 설계 대비 손실 열세, 회로의 불통과 사례)는 지우지 않고 설계 변경과 성립 조건의 근거로 남겼다. 남은 기간에는 16장의 검증 과제를 먼저 끝내고, 그 결과까지 넣어 최종 보고서로 정리한다.
